\documentclass[aspectratio=169]{beamer}
\usepackage{tikz}
\usepackage{graphicx}
\usepackage{listings}
\usepackage{xcolor}
\usepackage{hyperref}
\usepackage[spanish,es-noquoting]{babel}
\usepackage[round]{natbib}

\lstset{
  basicstyle=\ttfamily\small,
  keywordstyle=\color{blue},
  stringstyle=\color{red},
  commentstyle=\color{gray},
  showstringspaces=false
}
% Minimal look
\usetheme{Warsaw}
%\setbeamertemplate{headline}
\setbeamertemplate{footline}{
  \leavevmode%
  \hbox{%
    % Left: Section name
    \begin{beamercolorbox}[wd=.33\paperwidth,ht=2.25ex,dp=1ex,left]{author in head/foot}%
      \hspace*{2ex}\insertsection
    \end{beamercolorbox}%
    % Center: Presentation title
    \begin{beamercolorbox}[wd=.34\paperwidth,ht=2.25ex,dp=1ex,center]{title in head/foot}%
      \insertshorttitle
    \end{beamercolorbox}%
    % Right: Slide numbers
    \begin{beamercolorbox}[wd=.33\paperwidth,ht=2.25ex,dp=1ex,right]{date in head/foot}%
      \insertframenumber{} / \inserttotalframenumber\hspace*{2ex}
    \end{beamercolorbox}%
  }%
  \vskip0pt%
}

\usecolortheme{default}
\usefonttheme{serif}
\graphicspath{{../}{./}}
\setbeamertemplate{caption}[numbered]

% Title info
\title[Module 7: Information Retrieva]{Module 7: Information Retrieval}
\subtitle{Leveraging MS Learn via Copilot}
\author{Ricardo Reyes-Jimenez}
\institute{Capgemini}
\date{\today}



\begin{document}

% Title slide
\begin{frame}
  \titlepage
\end{frame}

% Reset background for normal slides
\usebackgroundtemplate{}



\begin{frame}[fragile]{Session Objective}


\textbf{By the end of this module, the Engagement Manager will be able to:}

1. Understand the limitations of standard Copilot for internal document retrieval.
2. Use Copilot to effectively query and summarize official \textbf{MS Learn} documentation.
3. Translate complex technical concepts for non-technical stakeholders.

\end{frame}



\begin{frame}[fragile]{1. Information Retrieval Reality for EMs}


\textbf{You have standard Copilot, not a custom Agent.}

\begin{itemize}
  \item \textbf{The Limitation:} You cannot train the model on thousands of internal Capgemini PDFs permanently.
  \item \textbf{The Workaround:} You must bring the context \textit{to} the chat, or direct Copilot to search the public web strategically (e.g., querying `learn.microsoft.com`) \cite{copilot2025}.
  \item \textbf{The Value:} EMs often need to justify architectural decisions or compliance standards to clients. MS Learn is the ultimate source of truth for Microsoft ecosystem projects.
\end{itemize}
\end{frame}



\begin{frame}[fragile]{Exercise 1: Policy Retrieval from MS Learn}


\textbf{Context:} The client is nervous about data governance in their new Azure deployment. You need the latest official Microsoft stance to reassure them.
\textbf{Task:} Prompt Copilot to search MS Learn and summarize the governance best practices.

\begin{itemize}
  \item \textbf{Draft Prompt to Try:}
\end{itemize}
\textit{``Search Microsoft Learn documentation for the latest best practices on 'Azure Data Governance'. Summarize the top 3 principles in a professional, reassuring email directed to a client's Chief Risk Officer. Include the source links at the bottom.''}

\end{frame}



\begin{frame}[fragile]{2. Bridging the Technical Gap}


\textbf{Engagement Managers are the bridge between developers and the business.}

\begin{itemize}
  \item Developers speak in acronyms; clients speak in ROI and Risk.
  \item Copilot excels at analogy generation and technical summarization, making you a more effective communicator.
\end{itemize}
\end{frame}



\begin{frame}[fragile]{Exercise 2: Technical Summarization via Analogy}


\textbf{Context:} Your technical lead insists on refactoring a monolithic application into microservices, but the client doesn't understand why it costs \$50,000 and delays the project.
\textbf{Task:} Use Copilot to explain the concept using a relatable business analogy.

\begin{itemize}
  \item \textbf{Draft Prompt to Try:}
\end{itemize}
\textit{``Explain the technical difference between a 'Monolithic Architecture' and 'Microservices'. Do not use heavy IT jargon. Instead, use an analogy related to managing a large corporate restaurant or kitchen. The goal is to justify to a non-technical CEO why migrating to microservices is worth the upfront investment.''}

\end{frame}



\begin{frame}[fragile]{Recap \& Next Session}


\textbf{Today we learned:}
\begin{itemize}
  \item How to direct Copilot to authoritative sources like MS Learn.
  \item How to use AI to translate dense technical architecture into business-friendly analogies.
\end{itemize}
\textbf{Next Session:}
\textit{Module 8: Final Integration (Role-Play).} The ultimate test. We will run a multi-step crisis simulation combining SOW analysis, Scope Creep detection, and diplomatic communication.
\end{frame}




\begin{frame}[allowframebreaks]{References}
  \bibliographystyle{apalike}
  \bibliography{references}
\end{frame}


\end{document}